
\chapter{Rolling Dice}



\section{Dice Rolls}
Players resolve the effects of character actions by rolling dice.
A \textit{Roll} is typically one die, but may involve multiple dice of the 
same and/or different types in a \textit{Dice Pool}.

\subsection{Natural Result}
The \textit{Natural Result} of a die roll is the value shown by the die 
after it settles and all Advantage is resolved.

\subsection{Critical Roll}
A \textit{Critical Roll} occurs when at least one die in a pool has a 
Natural Result that is the highest possible value for that type of die 
(e.g., a 6 on a \Roll{d6}, or a 12 on a \Roll{d12}).
Abilities or effects that trigger on a Critical Roll trigger once 
for each die in the pool that is a Critical Roll.

\subsection{Modifiers \& Final Result}
\textit{Modifiers} may be added to the Natural Result 
to increase or decrease the \textit{Final Result} of a roll.

\begin{DesignerNote}
  General modifiers (e.g., ``+1 to all rolls") are rare; 
  most modifiers impose a restriction on the type of die, 
  skill, or situation for the roll.
  For example, ``+1 to Clash rolls with Bladed weapons" is a modifier 
  that would only be applied when the character is rolling a \Roll{d6} 
  to resolve an action that uses their equipped Bladed weapon.
\end{DesignerNote}



\section{Advantage}
\textit{Advantage} is a mechanic that allows a player to build a dice pool 
for a roll, even when that roll is already a dice pool.

\vspace{0.4em}

By default, a roll's Advantage is 0 (Neutral).

\vspace{0.4em}

An ability that \textit{gains} Advantage increments this value, 
whereas one that \textit{loses} Advantage decrements it.
There may be multiple overlapping effects that cause a roll 
to gain or lose advantage;
the roll's overall Advantage is the sum of these effects.

\vspace{0.4em}

When a roll's Advantage is \textbf{greater than 0}, 
add that many dice of any type to the pool to be rolled,
then remove that many of the same type after rolling.

\begin{Example}
  If a roll has +1 Advantage for a dice pool of \Roll{1d6}, 
  instead the player rolls \Roll{2d6} and chooses one \Roll{d6}
  to remove and one to keep as the Natural Result.

  \vspace{0.4em}

  If a roll has +1 Advantage for a dice pool of \Roll{1d8}+\Roll{2d4},
  instead the player may roll \Roll{2d8}+\Roll{2d4} 
  or \Roll{1d8}+\Roll{3d4}, 
  then remove one \Roll{d8} or one \Roll{d4}, respectively.
\end{Example}

\vspace{0.4em}

When a roll's Advantage is \textbf{less than 0}, 
add that many more dice of the greatest type in the pool, 
then remove that many of the same type,
in highest order.

\begin{Example}
  If a roll has -1 Advantage for a dice pool of \Roll{1d6}, 
  instead the player rolls \Roll{2d6},
  then removes the \Roll{d6} with the highest value,
  keeping the remaining \Roll{d6} as the Natural Result.

  \vspace{0.4em}

  If a roll has -1 Advantage for a dice pool of \Roll{1d8}+\Roll{2d4},
  instead the player rolls \Roll{2d8}+\Roll{2d4} 
  then removes the \Roll{d8} with the highest value,
  keeping the remaining dice as the Natural Result.
\end{Example}



\section{Re-Rolling}
A \textit{Re-roll} allows the player to redo an entire roll sequence 
for all dice in that pool, including dice added by Advantage, 
applying effects and modifiers as if it were the original roll.

\vspace{0.4em}

If a rule only grants a re-roll for a certain type of die, 
re-roll only those dice as if they were their own pool,
then add them back into the original pool.

\vspace{0.4em}

If an effect would trigger as a result of the original roll, it is negated, 
and the triggering ability is evaluated for the new roll instead.

\vspace{0.4em}

Like Advantage, Re-roll can be stacked, 
making it a powerful mechanic that should be granted sparingly.



\section{Die Classes}
In W.I.Z.A.R.D.S., each type of die has a distinct role (pun intended).

\vspace{0.4em}

This design choice serves a twofold simplicity:\\
it is both easier for campaign module designers to balance game mechanics
and easier for players to remember which dice are required for their abilities.


\subsection{Clash (d6)}
The \textit{Clash Roll} is used when a character ability requires players to perform a contested roll.

\vspace{0.4em}

Every Clash follows a specific procedure:

\begin{enumerate}
  \item The player controlling the ability rolls \Roll{1d6} first. 
  If it is a Critical Roll, they automatically succeed, 
  and the defender does not get a saving roll.
  Otherwise, add the character's \textit{Clash} modifiers to the roll.
  \item Next, the player whose character is targeted by the ability 
  rolls \Roll{d6}.
  If it is a Critical Roll, they automatically save, regardless of modifiers.
  Otherwise, add the character's \textit{Clash Save} modifiers to the roll.
  \item If the attacker's roll is greater than the defender's roll, 
  they succeed.
  Otherwise, the defender saves their character. 
  (Tie goes to the defender.)
\end{enumerate}


\subsection{Damage/Restoration (d4/d8/d12)}
A \textit{Damage Roll} is always made by the player who controls 
the source of the damage, and a \textit{Restoration Roll} 
is always made by the player who controls the source of the restoration.

\vspace{0.4em}

Damage classes are standardized as \textit{Light} (\Roll{d4}), 
\textit{Normal} (\Roll{d8}), and \textit{Heavy} (\Roll{d12}) damage.

\vspace{0.4em}

A damaging ability or effect will specify the damage class 
and any type keywords associated with the damage,
and may deal more than one type or class of damage.
All damage die produced by a given effect are resolved simultaneously, 
as a single roll. 

\vspace{0.4em}

However, some abilities may have more than one effect, 
in which case damage is calculated separately.

\begin{Example}
  If a poisoned dagger produces Light (\Roll{d4}) damage from an attack, 
  and further Light (\Roll{d4}) damage from poison, 
  these are rolled separately, as the poison is its own effect.

  \vspace{0.4em}

  If a spiked club deals Heavy (\Roll{d12}) damage from the club, 
  combined with multiple instances of Light (\Roll{d4}) 
  damage from the spikes, these are resolved in the same roll as a pool.
\end{Example}

When dealing damage, identify the reservoir (e.g., Life or Mana) 
that is targeted by the damaging ability, 
then subtract the final result of the roll from that reservoir.
Unless otherwise specified, 
damage is assumed to target a character's Life reservoir.

\vspace{0.4em}

Follow the same procedure for restoration, 
but instead add the final result of the roll to the target reservoir.


\subsection{Outcome (d10/d100)}
An \textit{Outcome Roll} may be used to choose from a table of 
predetermined possibilities, 
or to determine the duration of a persistent effect.

\vspace{0.4em}

Tables may be used to represent the result of an effect, 
such as from a character ability or terrain hazard,
or it may represent the outcome of a character action, 
such as studying a particular artifact or environment.

\vspace{0.4em}

The \Roll{d10} is used for simple tables with limited outcomes, 
whereas \Roll{d100} may be used for longer lists,
distributions of outcome probability, 
and other percentile calculations of outcome.

\vspace{0.4em}

If a table is \textit{Cumulative}, 
then the character receives the effects of all values less than their roll.
The tables for cumulative rolls should be designated as such, 
otherwise rolls are assumed to be singular and specific to the value rolled.


\subsection{Story (d20)}
A \textit{Story Roll} is functionally a petition to a Game Master 
to allow a player's character to act outside of their explicit 
skills and abilities, in service of the plot.

\vspace{0.4em}

When a player describes the character's intended actions and outcome,
the GM considers the difficulty of the attempted action in the context 
of the situation,
then considers the suitability of the character's action 
to success in the intended outcome.

\vspace{0.4em}

Based on these considerations, 
the GM may arbitrarily assign a \textit{Threshold} 
that the Story roll must meet in order to succeed,
as well as an Attribute Modifier (see Chapter 4) 
that best applies to the roll.
Thresholds may be disclosed at the GM's discretion, 
and may cover any variety of outcomes along a spectrum of failure and success.

\vspace{0.4em}

The GM may also arbitrarily assign an Action cost (see Chapter 3: Combat)
to perform the Story action.

\vspace{0.4em}

If the terms are acceptable to the GM and to the player, 
the player may roll a \Roll{d20} to resolve the character's unique action.

